\section{Affine outputs and the original power-root chart}
\label{ai19-section}

This section gives complete ordinary intersection bounds for the
one-parameter chart of actual Eisenstein powers. The bounds include
integer residual coefficients. They concern a fixed exponent and this
root chart; they do not describe all two-coordinate roots.

Let \(\mathcal O=\mathbb Z[\zeta]\), where
\(\zeta^2-\zeta+1=0\) and \(\operatorname{Im}\zeta>0\). In the basis
\((1,\zeta)\), write \(\det\) for the coordinate determinant and
\(N(a+b\zeta)=a^2+ab+b^2\). Fix integers \(n\ge2\), \(B\ge n\), and set
\[
 Z(x)=(3x+\zeta)^n=A(x)+C(x)\zeta,\qquad x\ge B.
\]
All lines below are affine lines in these two real coordinates.

\begin{proposition}[Strict curvature of the complete chart]
\label{ai19-curvature}
Both \(A'(x)\) and \(C'(x)\) are positive, and
\[
 \det(Z'(x),Z''(x))
 =-27n^2(n-1)N(3x+\zeta)^{n-2}<0.
\]
Every real affine line meets this chart in at most two points. The same
bound holds after any invertible real linear change of coordinates.
\end{proposition}
\begin{proof}
Put \(\theta(x)=\arg(3x+\zeta)\). Then
\(0<\theta(x)<1/(3x)\), so \(0<(n-1)\theta(x)<\pi/3\).
Consequently \(Z'(x)=3n(3x+\zeta)^{n-1}\) has both coordinates positive.
Also \(Z''(x)=9n(n-1)(3x+\zeta)^{n-2}\).
Multiplication by an element \(W\) scales coordinate determinants by
\(N(W)\), while \(\det(3x+\zeta,1)=-1\). Factoring the common
\((3x+\zeta)^{n-2}\) proves the displayed determinant identity.

Since \(A'>0\), use \(A\) as a coordinate and write \(C=F(A)\). Then
\[
 F''(A(x))=\frac{\det(Z'(x),Z''(x))}{A'(x)^3}<0.
\]
A vertical line meets this graph at most once. Three intersections
with a nonvertical line would contradict strict concavity at the
middle abscissa, using the chord through the outer two intersections.
An invertible linear map takes affine lines to affine lines under its
inverse, so it preserves the bound.
\end{proof}

\begin{proposition}[Fixed-content and unfrozen affine selectors]
\label{ai19-selectors}
Let \(Z_i=Z(k_i)\), where \(B\le k_2<k_1\), put
\(D=Z_1-Z_2\) and \(\Delta=\det(Z_1,Z_2)>0\). In a fixed-content
selector progression
\[
 V_j=\frac{Z_2+(t_*+Hj)D}{g},\qquad j\ge1,
\]
with \(H,g>0\), the distinct outputs lie on the line
\[
 \det(D,V_j)=\Delta/g.
\]
At most two outputs equal \(Z(k)\) for an integer \(k\ge B\).
If \(g=1\) and the two nonempty disjoint CRT packets exclude the
selector parameters \(t=0,1\), there are no such equalities.

For all CRT selectors of these same two roots before content is frozen,
there are at most \(2\tau(|\Delta|)\) parameters whose primitive output
equals \(Z(k)\) for an integer \(k\ge B\).
\end{proposition}
\begin{proof}
Both coordinates of \(D\) are positive, and \(H/g>0\), so the
progression is injective. Taking its determinant with \(D\) gives
the stated line because \(\det(D,Z_2)=\Delta\).
Proposition~\ref{ai19-curvature} bounds the intersection by two.
For \(g=1\), the two known intersections are already \(Z(k_2)\)
and \(Z(k_1)\). No third intersection is possible, and equality to
these two outputs forces \(t=0\) and \(t=1\), respectively.

The actual content theorem for affine selectors,
Theorem~\ref{lm-exact-selector}, gives
\(g_t\mid|\Delta|\). For each possible positive divisor \(g_t\),
the outputs lie on its one determinant line and have at most two
possible values \(Z(k)\). For fixed content and output, the
nonzero vector \(D\) determines \(t\) uniquely. Sum over the
positive divisors of \(|\Delta|\). This is an upper bound only,
with no assertion that any of those intersections exists.
\end{proof}

\begin{theorem}[Uniform counting with every small integer residual]
\label{ai19-residuals}
Let \(V_j=V_0+jV_1\) be any affine integer progression in
\(\mathcal O\), with \(V_1\ne0\), and let \(X\ge1\) be real.
The number of positive integer parameters \(j\) admitting
\[
 V_j=\tau(3k+\zeta)^n,\qquad
 \tau\in\mathcal O\setminus\{0\},\quad N(\tau)\le X,\quad
 k\in\mathbb Z,\quad k\ge B
\]
is at most
\[
 2\left[\left(2\lfloor\sqrt{2X}\rfloor+1\right)^2-1\right]
 <30X.
\]
If the conjugate of \((3k+\zeta)^n\) is also allowed, the bound is
less than \(60X\). These constants are independent of \(n,B\),
the affine line, and any packet primes or depths. Unit multipliers
are already included among the possible \(\tau\).
\end{theorem}
\begin{proof}
For each fixed nonzero \(\tau\), multiplication by \(\tau\) is
invertible over \(\mathbb R\), with determinant \(N(\tau)>0\).
Its image of the power chart meets the progression line at most
twice by Proposition~\ref{ai19-curvature}. Since \(V_1\ne0\),
each output determines \(j\) uniquely. Summing these bounds may
overcount multiple representations, which is harmless.

For \(\tau=a+b\zeta\), one has
\(N(\tau)\ge(a^2+b^2)/2\). Thus both coordinates are at most
\(\lfloor\sqrt{2X}\rfloor\) in absolute value, giving at most
\((2\lfloor\sqrt{2X}\rfloor+1)^2-1\) nonzero pairs. For \(X\ge1\),
\[
 (2\sqrt{2X}+1)^2
 \le(9+4\sqrt2)X<15X.
\]
This proves the first bound. Conjugation is an additional invertible
real linear transformation, so summing the two orientation bounds
gives \(60X\). No power-free or coprimality restriction on the
residual was used; such restrictions could only reduce the count.
\end{proof}

\begin{corollary}[Scope of transfer from the paid affine family]
\label{ai19-transfer}
Apply Theorem~\ref{ai19-residuals} to a fixed-data progression in
Theorem~\ref{lm-paid-affine-family}. For every real \(L\ge0\), fewer
than \(60e^{nL}\) parameters in the entire progression admit either
orientation with
\[
 \log N(\tau)\le nL.
\]
More generally, if \(X(J)\ge1\) and \(X(J)=o(J)\), then only \(o(J)\)
indices \(1\le j\le J\) admit such a representation with
\(N(\tau)\le X(J)\). All but \(o(J)\) of the paid good indices
up to \(J\) have no such representation. In particular, for fixed
\(0<\eta<1\), the good indices with a representation satisfying
\(N(\tau)\le j^\eta\) form a zero-density subset.
\end{corollary}
\begin{proof}
Use \(X=e^{nL}\) for the first assertion and \(X=X(J)\) for the
second. Intersecting the resulting bound with the good set can
only decrease it; that good set has fixed-data density at least
\(1/4\) by Theorem~\ref{lm-paid-affine-family}. For the last assertion,
\(j\le J\) and \(N(\tau)\le j^\eta\) imply \(N(\tau)\le J^\eta\);
the exceptional count is therefore less than \(60J^\eta=o(J)\).
\end{proof}

These bounds explain a specific obstruction to identifying the
affine outputs with many same-index outputs in the original root
chart. They do not exclude a finite useful intersection, a different
two-coordinate root chart, a changing exponent, a transfer between
different outputs, or a nonlinear selector. The \(60X\) bound is
uniform, whereas the paid affine density has fixed-data thresholds.
No uniform statement for moving packets, no original far-tail
budget, and no control of the dependent-root complement follows.
The results in this section are ordinary proofs, not a claim of
complete Lean formalization.
